\section{Analysis protocol as predeclared}
\label{app:protocol}
\TODO{FOR BIODIVERSITY AND CONSERVATION: this appendix is not part of the submitted manuscript. It becomes a separate file, Online Resource 1, uploaded as \texttt{ESM\_1.pdf}. Springer asks that every supplementary file include, at the top, the article title, journal name, author names, and the corresponding author's affiliation and email, plus a short caption describing the file. In the main text it is referred to as ``Online Resource 1''; typing \texttt{\textbackslash suppA} prints that automatically in the submission layout and ``Appendix A'' in the capstone.}
\TODO{PLAIN-LANGUAGE GUIDE FOR THE PROTOCOL APPENDIX.
The saved protocol is tagged \texttt{protocol-v1.0} in the
\texttt{cheetah-connectivity} repository at commit
\texttt{d7de8675a074dfbffc0c7c856bb5645e1e13867d}, dated 18 August 2026. However, the
file itself still says version 0.1 draft. Do not change the old file. Explain that the
Git tag marked the saved version even though its internal label was not updated.
Then list, in simple terms, what changed between the plan and the final analysis: the
core definition, sensitivity tests, definition of a stable link, scenario names,
network measurements, protected-area comparison, monitoring index
that was not built, and 2030 scenarios that were added later. Include the objective
and hypothesis outcomes: O1 was completed; O2 was narrowed; O3 was replaced rather
than completed; H1 was mixed; H2 did not pass its full rule; H3 was only partly tested
and did not match the ``limited subset'' prediction; and H4 was not tested.}

\TODO{ADDED 5 OCTOBER 2026: FOUR MORE ITEMS FOR THE CHANGE LIST, plus one correction.\\
1. A MAIN SURFACE WAS CHOSEN. The protocol said no resistance scenario would be treated as the main answer. The paper uses the vegetation-balanced surface with documented fences as the main one and reports the rest as sensitivity. Say this was a change and give your reason in one sentence.\\
2. WHEN THE WEIGHTS WERE SET. Say plainly whether the final weights (60/20/20 and 30/30/30/10) were fixed before or after you had looked at resistance maps or results. This is the first thing a reader checks in a predeclared study. Either answer is acceptable if it is stated honestly; if any weight was adjusted after seeing outputs, say which one and why.\\
3. PC/dPC WAS DROPPED. The protocol planned probability of connectivity at 50, 100, 200, and 400~km. Only betweenness was calculated. Give the real reason (see the TODO in the Network structure subsection).\\
4. ONE CORE DEFINITION, NOT THREE. The protocol planned three core definitions with minimum areas of 2,000, 1,000, and 500~km$^2$. The analysis used one density-based definition. List it here if it is not already covered by ``core definition'' above.\\
CORRECTION: the old list included ``occurrence check'' as a change. That is no longer true. The planned independent tracking check (protocol steps 23.1--23.5) was completed using the Limpopo data. What is still different from the plan: the Namibia camera-trap comparison was not done, and two kinds of analysis were added. First, a second, same-archive occurrence check. Second, four stricter follow-ups to the tracking check (a step check comparing each real move with moves the animal could have made instead, a breakdown by input layer, a repeat on all nine surfaces, and a road-crossing count). List these as additions, say they were added after the planned check was run, and say that none of them changed any weight or the choice of main surface.}
